\documentclass[11pt,a4paper]{article}
\usepackage[margin=27mm]{geometry}
\usepackage{amsmath,amssymb,amsthm,fontspec,unicode-math,hyperref,booktabs,longtable}
\usepackage{array,calc,fvextra}
\setmainfont{STIXGeneral}
\setmathfont{STIX Two Math}
\setmonofont{Menlo}
\hypersetup{hidelinks}
\setlength{\emergencystretch}{2em}
\providecommand{\tightlist}{\setlength{\itemsep}{0pt}\setlength{\parskip}{0pt}}
\newtheorem{theorem}{Theorem}
\newtheorem{lemma}[theorem]{Lemma}
\newtheorem{corollary}[theorem]{Corollary}
\newtheorem{proposition}[theorem]{Proposition}
\newcommand{\gen}[1]{\langle #1\rangle}
\newcommand{\Sym}{\operatorname{Sym}}
\newcommand{\Aut}{\operatorname{Aut}}
\newcommand{\Fix}{\operatorname{Fix}}

\title{Recognizable extensions of colored graphs}
\author{Human Compression Project}
\date{}
\begin{document}
\maketitle
\begin{abstract}
A construction can carry every symmetry of its input and still acquire additional
symmetries. We separate these two questions by a restriction and extension
criterion for permutation groups and their distinguished subgroups. Intrinsic
recognition of retained vertices, rigidity over those vertices, and explicit
extension identify the full group. Component boundaries provide a useful local
rigidity test. We give complete proofs for the chamber completion and cleaning
constructions of D\'iaz Ramos, Molinier and Viruel. The first application recovers their finite normal-pair representation theorem
and recursive symmetric-group family. A second application treats the
constructions of Leemans, Stokes and Tranchida. We include the full kernel
created by repeated shadows, and the free-basis coset geometry and normalizer
arguments needed by its infinite examples. Distinct-length arms preserve
ordered attachments, while a thin-doubling construction separates the
connected-chamber splitting theorem from the unqualified statement.
\end{abstract}
The first four sections concern finite incidence systems. The free-basis
sections explicitly allow infinitely many objects with a finite palette;
the finite convention in the opening definitions is local to the finite
constructions. The thin-doubling proposition returns to finite systems.
\section{Restriction, extension and component boundaries}
\label{sec:recognition}
An incidence system in this paper is a nonempty finite simple graph $X$, together
with a surjective proper coloring $t:V(X)\to I$. Thus adjacent vertices have
different colors. A \emph{correlation} is a graph automorphism $f$ that permutes
the color classes: there is a permutation $\sigma_f$ of $I$ such that
$t(f(v))=\sigma_f(t(v))$ for all $v$. Surjectivity makes $\sigma_f$ unique.
Write $C(X)$ for the correlation group and $K(X)$ for its subgroup preserving
each color. The map $f\mapsto\sigma_f$ is a homomorphism, so $K(X)$ is normal
and $C(X)/K(X)$ embeds in $\Sym(I)$. In general $C(X)$ need not be the full
uncolored graph automorphism group.

A \emph{chamber} is a clique meeting every color class. An incidence system is
an \emph{incidence geometry} when each clique, including the empty clique, is
contained in a chamber. No transitivity, thinness or residual connectedness is
part of this definition. These are the conventions of
D\'iaz Ramos--Molinier--Viruel~\cite{DMV}, Definitions 2.1--2.3.

The following elementary fact is the common group-theoretic argument. It is
stated for arbitrary permutation groups; finiteness enters the graph
constructions later.

\begin{theorem}[Restriction and extension]\label{UR.Restriction}
Let $\iota:V\hookrightarrow W$ be an injection, let
$A\leq\Sym(V)$ and $B\leq\Sym(W)$, and let $A_0\leq A$ and $B_0\leq B$.
Suppose that:
\begin{enumerate}
\item for each $b\in B$ there is $a\in A$ such that
$b\iota=\iota a$;
\item a member of $B$ fixing $\iota(V)$ pointwise is the identity;
\item every $a\in A$ has an extension $b\in B$ satisfying
$b\iota=\iota a$;
\item for such a pair $(b,a)$, one has $b\in B_0$ if and only if $a\in A_0$.
\end{enumerate}
Then restriction induces an isomorphism $B\cong A$ carrying $B_0$ onto $A_0$.
Every extension in the third condition is unique, and the resulting extension
map is the inverse homomorphism.
\end{theorem}
\begin{proof}
The injection $\iota$ makes the element $a$ in the first condition unique;
write $\rho(b)=a$. For $b,c\in B$,
\[
 bc\iota=b\iota\rho(c)=\iota\rho(b)\rho(c),
\]
so uniqueness gives $\rho(bc)=\rho(b)\rho(c)$. The same argument gives
$\rho(1)=1$. Its kernel consists exactly of the elements fixing $\iota(V)$
pointwise, and is trivial by the second condition. The third condition makes
$\rho$ surjective. The last condition gives
$\rho(B_0)=A_0$: containment follows immediately, while a preimage of any
$a\in A_0$ belongs to $B_0$. Thus $\rho$ is the asserted pair isomorphism.
Injectivity gives uniqueness of extensions, and the inverse of an isomorphism
is a homomorphism.
\end{proof}

For colored graphs $X,Y$, one applies this theorem to $A=C(X)$,
$B=C(Y)$ and their palette kernels. Recognizing a retained copy of $V(X)$
is not enough: the restricted permutation must preserve the \emph{reconstructed}
old incidences and colors. In particular, old edges need not remain edges of
$Y$. Conversely, an explicit action of $C(X)$ on $Y$ establishes only the
extension condition, until rigidity rules out further correlations.

Here is a local criterion for that rigidity. All induced graphs and boundaries
in its statement are taken in the constructed graph $Y$.

\begin{lemma}[Recognition by component boundaries]\label{UR.Components}
Let $S\subseteq V(Y)$ be invariant under a group $B$ of correlations of $Y$.
For each connected component $D$ of $Y\setminus S$, define its boundary
\[
 \partial D=\{s\in S:\text{$s$ has a neighbor in $D$}\}.
\]
Assume distinct components have distinct boundaries. Suppose also that any
correlation of the induced colored graph on $D\cup\partial D$ fixing
$\partial D$ pointwise fixes $D$ pointwise. Then a member of $B$ fixing $S$
pointwise is the identity on $Y$.
\end{lemma}
\begin{proof}
Such a member permutes the components of $Y\setminus S$ and sends each
boundary to itself. Distinctness of the boundary sets therefore fixes each
component setwise. Its restriction to $D\cup\partial D$ is a correlation,
fixes the boundary pointwise, and is the identity by the local assumption.
This holds for every component; it is already the identity on $S$.
\end{proof}

Only color classes occurring in an induced subgraph are used when speaking of
its restricted coloring. A global correlation preserving that subgraph does
permute those classes, even if it permutes other colors outside the subgraph.
Thus the local hypothesis in Lemma~\ref{UR.Components} does not silently
replace correlations by color-preserving automorphisms.

\section{Completing edges to chambers}
\label{sec:chambers}
The following construction and theorem are due to D\'iaz Ramos, Molinier and
Viruel~\cite[Theorem 4.1]{DMV}. We retain their construction and prove all four
conditions of Theorem~\ref{UR.Restriction}.

\begin{theorem}[Chamber completion]\label{UR.Chambers}
Let $X$ be a finite properly colored triangle-free graph of minimum degree at
least two, with palette $I$ of size $r$. For every edge $e=\{v,w\}$ and every
$i\in I\setminus\{t(v),t(w)\}$, introduce a vertex $u_{e,i}$ of color $i$.
Make
\[
 B_e=\{v,w\}\cup\{u_{e,i}:i\notin\{t(v),t(w)\}\}
\]
a complete graph, retaining all old vertices and adding no other edges.
The resulting graph $J(X)$ is an incidence geometry on the same palette and
restriction gives
\[
 (C(J(X)),K(J(X)))\cong(C(X),K(X)).
\]
It has $|V(X)|+(r-2)|E(X)|$ vertices and the same connected components as $X$.
\end{theorem}
\begin{proof}
Every $B_e$ has exactly one vertex of each color. A clique containing a new
vertex lies entirely in its $B_e$, since that vertex has no neighbors outside
$B_e$. A clique consisting of old vertices has size at most two by
triangle-freeness. A two-vertex clique supplies its edge $e$; a singleton
extends using an incident edge; the empty clique extends because $X$ has an
edge. Hence every clique is contained in a chamber.

A new vertex has degree $r-1$. At an old vertex $v$, the sets
$B_e\setminus\{v\}$, for edges $e$ through $v$, are pairwise disjoint.
Consequently
\begin{equation}\label{eq:degrees}
 \deg_{J(X)}(v)=\deg_X(v)(r-1).
\end{equation}
The hypotheses imply $r\geq2$ and $\deg_X(v)\geq2$, so the old vertices
are exactly those of degree greater than $r-1$. Every correlation recognizes
them. Their induced graph is $X$, and their old colors are unchanged; the
restriction therefore belongs to $C(X)$.

Suppose a correlation fixes every old vertex. Every old color occurs, so it
fixes every color. If a new vertex has old neighbor set $\{v,w\}$ and color
$i$, it must be $u_{\{v,w\},i}$. Thus all new vertices are fixed and
restriction is injective. If $f\in C(X)$ induces $\sigma_f\in\Sym(I)$,
define
\[
 \widetilde f(v)=f(v),\qquad
 \widetilde f(u_{e,i})=u_{f(e),\sigma_f(i)}.
\]
This is a bijection, with inverse defined from $f^{-1}$. It carries each
$B_e$ onto $B_{f(e)}$, and hence preserves adjacency in both directions.
It fixes all colors precisely when $f$ does. Theorem~\ref{UR.Restriction}
now identifies the pairs. When $r=2$ no vertices are added, and this proof
still applies.

There are $r-2$ new vertices per edge, giving the count. Every new vertex is
attached within the component of its edge, and all old edges are retained.
No components are joined or lost.
\end{proof}

The source's degree argument on physical page 7 uses $r$ in place of $r-1$.
Equation~\eqref{eq:degrees} is the actual degree calculation. The strict
separation needed for its theorem is unaffected. The statement concerns this
construction and these hypotheses; it does not assert that they are logically
necessary for every conceivable pair-preserving completion.

\section{Cleaning arbitrary incidence systems}
\label{sec:cleaning}
The next construction is \cite[Theorem 4.2]{DMV}. We include the rigidity of
its tail and the isolated-vertex case, since those are the local obligations
that make its symmetry assertion valid.

\begin{theorem}[Cleaning]\label{UR.Cleaning}
Every nonempty finite properly colored graph $X$ has a finite properly colored,
triangle-free graph $L(X)$ of minimum degree at least two satisfying
\[
 (C(L(X)),K(L(X)))\cong(C(X),K(X)).
\]
It preserves the connected-component decomposition and adds exactly
$q(2M+5)$ vertices, where
\[
 M=\max_{i\in I}|t^{-1}(i)|,
 \qquad W=\{v:\deg_X(v)\leq1\},\qquad q=|E(X)|+|W|.
\]
Its palette consists of the old colors and two new colors.
\end{theorem}
\begin{proof}
Since $X$ is nonempty, $M\geq1$; also $q\geq1$, since otherwise a nonempty
edgeless graph would have no vertices of degree at most one.
For each edge $e$ and each $w\in W$ take a separate gadget with vertices
$u_0,\ldots,u_{2M+4}$, consecutive edges $u_ju_{j+1}$, and the two chords
\begin{equation}\label{eq:tail}
 u_{2M+4}u_{2M-1},\qquad u_{2M+3}u_{2M}.
\end{equation}
Delete all old edges. For $e=\{v,w\}$ join both endpoints to the root
$u_{e,0}$ of its edge gadget. Join each $w\in W$ to $u_{w,0}$; if $w$ was
isolated, also join it to $u_{w,2}$. Retain old colors, color even gadget
indices by $\alpha$ and odd indices by $\beta$, where $\alpha,\beta$ are
distinct new colors.

Every new edge joins different colors. An old vertex of degree at least two
retains that degree; a former leaf acquires a second neighbor; a former
isolated vertex acquires its two displayed neighbors. Every gadget vertex has
degree at least two. There are no old--old edges. The new induced graph is
bipartite. A triangle containing an old vertex would require two adjacent new
neighbors in one gadget, but the only old vertex with two neighbors in one
gadget was isolated, and its neighbors $u_0,u_2$ are not adjacent. Thus $L(X)$
is triangle-free and has the required minimum degree.

The two new color classes have sizes $q(M+3)$ and $q(M+2)$. These sizes are
distinct and both exceed $M$, the largest old class size. Every correlation
therefore fixes each of these classes and preserves the old vertex set. Two
distinct old vertices were adjacent in $X$ exactly when they have a common
new neighbor in $L(X)$: their edge gadget supplies its root, and every other
gadget has only one old neighbor. This reconstructs the old graph as well as
its color partition, so restriction belongs to $C(X)$.

To establish rigidity, delete all old vertices. The components are precisely
the separate gadgets. Each edge gadget has its old endpoint pair as boundary;
each additional gadget has its one old vertex as boundary. The graph is simple
and at most one additional gadget is attached to a given old vertex, so all
these boundaries are distinct.

Consider a correlation of a gadget together with its boundary that fixes the
boundary. For an edge gadget or a former leaf, the root is the unique gadget
vertex adjacent to the boundary. For an isolated vertex the two attached
vertices $u_0,u_2$ have different degrees in this induced graph: the first has
degree two and the second at least three. Hence the root is fixed in all cases.
Starting at the root forces successive vertices until $u_{2M-1}$. At that
vertex the forward neighbors $u_{2M}$ and $u_{2M+4}$ have induced-gadget
degrees three and two respectively, so both are fixed. At $u_{2M}$ the
remaining neighbors $u_{2M+1}$ and $u_{2M+3}$ have induced-gadget degrees two
and three respectively. Finally $u_{2M+2}$ is forced. These branch degrees
are taken within the gadget, after its component has been recognized; an
isolated attachment therefore causes no exception when $M=1$.
Lemma~\ref{UR.Components} now gives pointwise rigidity over all old vertices.

A correlation $f$ of $X$ preserves degrees. Extend it by
\[
 u_{e,j}\longmapsto u_{f(e),j},\qquad
 u_{w,j}\longmapsto u_{f(w),j}.
\]
The complete edge list, including the isolated attachment, is preserved, and
the inverse is the same construction from $f^{-1}$. Both new colors are fixed;
old colors are fixed exactly when $f$ fixes them. Thus
Theorem~\ref{UR.Restriction} gives the asserted pair isomorphism.
There are $q$ gadgets of $2M+5$ vertices. Replacing an old edge preserves
connectivity between its endpoints and does not join different old components;
all other gadgets attach to one old vertex. This proves the component claim.
\end{proof}

\begin{corollary}\label{UR.Conversion}
Every finite incidence system has a finite incidence geometry with the same
correlation group and the same distinguished type-preserving subgroup.
Connected inputs give connected outputs.
\end{corollary}
\begin{proof}
Apply Theorem~\ref{UR.Chambers} to $L(X)$ from Theorem~\ref{UR.Cleaning}.
Both restriction maps identify the distinguished kernels. Each construction
preserves the component decomposition.
\end{proof}

\section{Repeated shadows and the restriction kernel}
\label{sec:shadows}
This section includes the elementary shadow-fiber theorem of the earlier
phase196 manuscript. It makes explicit the rigidity condition in
Theorem~\ref{UR.Restriction}; the old objects being recognizable does not
force the new objects to be determined by them.

Let $P$ be a finite typed incidence system, $J$ a finite set of new types,
and $Q_j$ pairwise disjoint finite sets, disjoint also from $P$. For
$q\in Q_j$ prescribe its \emph{complete old shadow} $s_j(q)\subseteq P$.
Retain the old incidences, put $p*q$ exactly when $p\in s_j(q)$, and make
all objects of distinct new types incident. There are no incidences within
a new type. Write $\Gamma$ for the extension and
$m_j(S)=|s_j^{-1}(S)|$. Let
\[
H=\{h\in K(P):m_j(hS)=m_j(S)\text{ for every }j,S\}.
\]
Absent shadows have multiplicity zero.

\begin{theorem}[The complete shadow kernel]\label{UR.Shadows}
Restriction to $P$ gives the split exact sequence
\[
1\longrightarrow\prod_{j,S}\Sym(s_j^{-1}(S))
\longrightarrow K(\Gamma)\longrightarrow H\longrightarrow1.
\]
Consequently restriction is faithful exactly when every nonempty shadow
fiber has one object. A splitting depends on chosen enumerations of the
fibers; it is not asserted to be canonical.
\end{theorem}
\begin{proof}
A type-preserving correlation preserves $P$, since old and new types are
disjoint. If its restriction is $h$, incidence preservation says that it
maps the fiber over $(j,S)$ bijectively to that over $(j,hS)$. Thus $h\in H$.
When $h=1$, it preserves each fiber. Conversely arbitrary independent
permutations of the fibers preserve old--new incidence and all new--new
incidences. This proves precisely the asserted kernel, not merely a
subgroup contained in it.

For each nonempty fiber choose an enumeration
$e_{j,S}:\{1,\ldots,m_j(S)\}\to s_j^{-1}(S)$. Extend $h\in H$ by
\[
 e_{j,S}(a)\longmapsto e_{j,hS}(a).
\]
Equality of multiplicities makes this well-defined. It preserves
old--new incidence because $p\in S$ if and only if $h(p)\in hS$;
the other incidences are preserved by construction. Its inverse is the
extension of $h^{-1}$. Applying successive extensions keeps the index $a$
unchanged, so these extensions form a homomorphism and split restriction.
Finally a symmetric group is trivial exactly when its set has at most one
element.
\end{proof}

For full correlations one must separately recognize the old object set
and allow any actual permutation of new types. Thus a signature that
includes a new type does not establish pointwise rigidity if a correlation
fixing old objects could still permute those types. The finite examples
in the second application verify this distinction explicitly.

\begin{corollary}[Recognition by unique shadows]\label{UR.UniqueShadows}
Suppose every correlation of an extension $Y$ recognizes an old structure
$X$, and its restriction preserves all reconstructed old incidences and
colors. Suppose every new object is uniquely determined by its complete
old shadow, even across any new types that correlations may permute.
Let $G\leq C(X)$ be exactly the old actions whose forced permutations of
new objects preserve the prescribed incidences and induce a permutation
of the full palette. Then restriction identifies $C(Y)$ with $G$; $K(Y)$
corresponds to the kernel of the resulting palette action.
\end{corollary}
\begin{proof}
Unique shadows give pointwise rigidity, and the definition of $G$ gives
exactly the permissible extensions. Theorem~\ref{UR.Restriction} applies.
\end{proof}

\begin{proposition}[Keeping ordered positions]\label{UR.OrderedShadows}
Let $P$ be a typed incidence system and $T$ a set of distinct ordered
$d$-tuples of old objects, with a prescribed type for every tuple. Form a
new system with a center for each tuple $t$ and a port $(t,i)$ for each
position $1\leq i\leq d$. Give ports distinct position types, disjoint
from the old and center types, and join $(t,i)$ just to its center and the
old object $t_i$. Keep the old incidences. In the subgroup preserving the
old types, center types and each position type, restriction is an
isomorphism onto the old automorphisms preserving the tuple sets by type.
\end{proposition}
\begin{proof}
For a correlation in the stated subgroup, the unique center and old
neighbor of a port show that the center image has entries $h(t_i)$ in
all positions. Since tuples are distinct, its image and all its ports
are forced. Conversely an old automorphism preserving the tuple sets
extends by $t\mapsto(h(t_1),\ldots,h(t_d))$ and $(t,i)\mapsto(h(t),i)$.
Both maps preserve all declared incidences and are inverse under
restriction. Apply Theorem~\ref{UR.Restriction}.
\end{proof}
This is a precise repair of loss of order at the level of a specified
role-preserving subgroup. It does not assert that the entire correlation
group recognizes those roles, that projective dualities extend, or that
the resulting incidence system is a geometry. Those are additional
requirements, not consequences of this proposition.

\begin{theorem}[Rigid arms encode ordered attachments]\label{UR.RigidArms}
Let $X$ be a finite properly colored graph and let $T$ be a family of
pairwise distinct ordered $d$-tuples of its vertices, with $d\ge3$ and
with a color $c(t)$ assigned to each tuple. Add a center for each tuple
and join it to its $i$th entry by a fresh path of length $\ell_i\ge2$,
where the $\ell_i$ are pairwise different. Give centers their prescribed
colors and give internal path vertices separate colors indexed by
$(i,k)$, their arm and distance from the center. All new colors are
disjoint from the old colors. No other new incidences are added.
Suppose every correlation of the resulting graph $Y$ preserves the old
vertex set $V(X)$. Then restriction identifies $C(Y)$ with the subgroup
of $C(X)$ preserving $T$ and its partition into tuple colors. Under this
identification $K(Y)$ is the kernel of the action on the old colors and
tuple colors; all internal path colors are fixed automatically.
\end{theorem}
\begin{proof}
Each component after deleting the old vertices is a subdivided star with
one vertex of degree $d$, its center. Its arm lengths are
$\ell_1-1,\ldots,\ell_d-1$ and are pairwise different. Every isomorphism
between two such components sends the center to the center, preserves
arm indices by their lengths and preserves each distance from the
center. Looking also at the unique old endpoint of each arm therefore
recovers the entire ordered tuple, not merely its unordered boundary.
Distinct tuples have distinct ordered boundaries, so a correlation
fixing all old vertices fixes each component pointwise.

A correlation restricts to $C(X)$ because all old incidences and colors
are retained. Its action sends each tuple to the ordered tuple of its
endpoint images, with one common permutation of the tuple colors.
Conversely an old correlation satisfying these conditions extends by
that rule on centers and by matching indices and distances on arms.
This is an incidence-preserving bijection, and the inverse is obtained
from the inverse old correlation. The internal colors are fixed and
the other colors transform just as asserted. Apply
Theorem~\ref{UR.Restriction}.
\end{proof}

\section{Free-basis incidence
systems}\label{free-basis-incidence-systems}

The following argument reuses the complete free-basis proof from the
earlier phase197 manuscript; its reduced-word and tree arguments are
included here in full. Classical reduced-word normal forms and the
universal property of free groups remain explicit inputs. No historical
priority is asserted. The object set may now be infinite, but the type
set remains finite.

Let \(I\) be a finite set of cardinality \(r\ge2\), let
\(F=F(s_i:i\in I)\) be the free group on this basis, and define \[
F_i=\langle s_j:j\ne i\rangle,
\qquad F_J=\bigcap_{j\in J}F_j.
\] The objects of type \(i\) are the right cosets \(F_i g\). Objects of
distinct types are incident when the cosets intersect. Denote this
incidence system by \(\Gamma_I\). Use the left permutation convention
for the natural action: \[
R_g(F_i h)=F_i h g^{-1}.
\tag{F1}
\] Thus \(g\mapsto R_g\) is a homomorphism.

\section{Free-basis cosets and the Helly
condition}\label{free-basis-cosets-and-the-helly-condition}

A standard free factor generated by a subset of the basis consists
exactly of the elements whose reduced words use only that subset. In
particular \[
F_J=\langle s_k:k\in I\setminus J\rangle.
\tag{F2}
\] This includes \(F_\varnothing=F\) and \(F_I=1\).

For a reduced word \(w\) and \(i\in I\), let \(P_i(w)\) be the prefix
ending at its last occurrence of \(s_i\) or \(s_i^{-1}\); set \(P_i(w)\)
to the empty word if there is no such occurrence.

\textbf{Lemma R2.Prefix.} For \(i\notin J\), \[
w\in F_JF_i
\quad\Longleftrightarrow\quad
P_i(w)\text{ uses no letter indexed by }J.
\tag{F3}
\]

\textbf{Proof.} If \(w=ab\) with \(a\in F_J\) and \(b\in F_i\), reduce
the concatenation of their reduced words. Cancellation takes place
across their common boundary. The remaining word is a prefix of the
reduced word of \(a\) followed by a suffix of that of \(b\), with no
further cancellation. The second part has no occurrence indexed by
\(i\), so the prefix ending at the last such occurrence lies entirely in
the first part. It avoids \(J\). This also covers complete cancellation
of the first part, when the prefix is empty. Conversely, if the
specified prefix avoids \(J\), split \(w\) immediately after its last
\(i\)-occurrence. The prefix belongs to \(F_J\), and the suffix has no
\(i\)-occurrence and belongs to \(F_i\). \(\square\)

\textbf{Corollary R2.Product.} For every \(J\subseteq I\) and
\(i\notin J\), \[
F_JF_i=\bigcap_{j\in J}F_jF_i.
\tag{F4}
\]

\textbf{Proof.} By (F3), membership in the right side says that
\(P_i(w)\) avoids each \(j\in J\), which says exactly that it avoids
\(J\). The empty intersection causes no exception. \(\square\)

We prove flag-transitivity directly, so no outside coset criterion is a
premise here. Consider the Cayley tree with edges joining \(g\) to
\(s_i^{\pm1}g\). Every right coset of a standard free factor is a convex
subtree, because the unique reduced path between its vertices uses only
its allowed generators. A finite family of pairwise intersecting
subtrees of a tree has a common vertex. For example, root the tree and
select, among the nearest-to-root vertices of the subtrees, one at
maximal depth. For any other subtree, take a point in its intersection
with the selected subtree. Both nearest vertices lie on the root path to
that point; the deeper one lies in both subtrees. It therefore lies in
every member of the family. This proves that every flag of \(\Gamma_I\)
consists of cosets with a common representative. Right translation sends
that representative to the identity and the flag to the standard flag of
its type. In particular every flag extends to a chamber.

\textbf{Theorem R2.Geometry.} The incidence system \(\Gamma_I\) is a
flag-transitive, residually connected geometry. The natural group
\(R(F)\) acts regularly on chambers.

\textbf{Proof.} The common-representative argument proves geometry and
flag-transitivity. The stabilizer of the standard chamber is
\(\bigcap_i F_i=1\) by (F2). Thus the chamber action is regular and the
natural action is faithful.

Translate an arbitrary flag to the standard flag of some type \(J\). Its
chambers are uniquely indexed by \(F_J\), and its remaining objects
identify with the cosets of \[
F_J\cap F_i=F_{J\cup\{i\}},\qquad i\notin J,
\] inside \(F_J\). To verify this identification, intersect each
incident coset with \(F_J\); it is nonempty by the common-representative
property and is a coset of the displayed intersection subgroup. Distinct
right cosets remain distinct after nonempty intersection.

The residue is consequently the same free-basis construction on the
remaining basis \(\{s_i:i\notin J\}\). If at least two types remain, its
chamber graph is connected: chambers labeled \(h\) and \(s_i^m h\) share
all objects except possibly type \(i\), and the remaining basis elements
generate \(F_J\). Each chamber is a clique with at least two vertices,
and consecutive chambers have a common object. The union of this
connected family of cliques is the entire residue incidence graph. Thus
every residue of rank at least two is connected. \(\square\)

This proof also yields exactly the generation condition quoted as
Theorem 7.4 in the source: \[
F_J=\left\langle F_{J\cup\{i\}}:i\in I\setminus J\right\rangle
\quad\text{when }|I\setminus J|\ge2.
\tag{F5}
\] Each remaining basis letter belongs to one of those factors by
choosing a different remaining index.

\textbf{Proof correction.} Membership in \(F_jF_i\) does not imply that
the entire reduced word avoids \(s_j\). Its second factor is allowed to
contain that letter. Theorem 7.3's printed proof uses this implication.
Formula (F3), rather than avoidance by the whole word, supplies the
valid common argument.

\section{Palette actions of a free
basis}\label{palette-actions-of-a-free-basis}

Let \(C=C(\Gamma_I)\) and \(K=K(\Gamma_I)\) denote the full correlation
and type-preserving groups. An automorphism \(\alpha\) of \(F\)
satisfying \(\alpha(F_i)=F_{\sigma(i)}\) induces the correlation \[
F_i g\longmapsto F_{\sigma(i)}\alpha(g).
\tag{F6}
\] The formula is well-defined because equality of cosets is preserved
by \(\alpha\), and incidence is preserved because \(\alpha\) preserves
nonempty intersection.

\textbf{Theorem R2.TypeImage.} The action of the full correlation group
on the palette has image \(\operatorname{Sym}(I)\). Thus \[
C/K\cong\operatorname{Sym}(I).
\tag{F7}
\]

\textbf{Proof.} Every permutation \(\sigma\) of the free basis extends
to an automorphism of \(F\) by the universal property of the free group,
and it sends \(F_i\) to \(F_{\sigma(i)}\). Formula (F6) therefore
realizes every palette permutation. The kernel of the palette action is
precisely \(K\) by definition. \(\square\)

There is no claim here that (F7) determines \(K\) or the entire group
\(C\). The natural subgroup \(R(F)\) need not be the full
type-preserving group.

The inherited Lean file \texttt{BasisPermutation.lean} verifies the
generic basis-permutation automorphism, its faithful dependence on the
permutation, and the equation transporting the subgroup closure
generated by every complementary basis. This is the formal
group-theoretic step in the proof of (F7). The coset quotient
construction and the inference of the full type quotient are written
proofs, not part of that file's formal scope.

\section{The normalizer of the chamber-regular
group}\label{the-normalizer-of-the-chamber-regular-group}

Write \[
W_I=C_2^I\rtimes\operatorname{Sym}(I)
\] for the signed basis permutations. Each maps \(s_i\) to
\(s_{\sigma(i)}^{\epsilon_i}\) with \(\epsilon_i\in\{1,-1\}\) and acts
by (F6).

\textbf{Theorem R2.Normalizer.} The normalizer of the natural
chamber-regular group is \[
N_C(R(F))=R(F)\rtimes W_I.
\tag{F8}
\] The signed basis subgroup is precisely the stabilizer of the standard
chamber inside this normalizer.

\textbf{Proof.} If \(c\) normalizes \(R(F)\), compose \(c\) with the
unique natural translation taking its image of the standard chamber back
to that chamber. It suffices to consider a normalizing correlation \(d\)
fixing the standard chamber. Conjugation by \(d\) induces an
automorphism \(\alpha\) of \(F\), since \(R:F\to R(F)\) is an
isomorphism. Let \(\sigma\) be the palette permutation of \(d\).

The pointwise stabilizer in \(R(F)\) of the standard chamber with type
\(i\) deleted is \[
R\!\left(\bigcap_{j\ne i}F_j\right)=R(\langle s_i\rangle).
\] Since \(d\) fixes the chamber as a set and permutes its types by
\(\sigma\), it sends this stabilizer to
\(R(\langle s_{\sigma(i)}\rangle)\). Thus \(\alpha\) maps each infinite
cyclic subgroup \(\langle s_i\rangle\) onto
\(\langle s_{\sigma(i)}\rangle\). An isomorphism of infinite cyclic
groups sends a generator to a generator, so
\(\alpha(s_i)=s_{\sigma(i)}^{\pm1}\). Hence \(\alpha\in W_I\).

A normalizing correlation fixing the standard chamber is determined by
\(\alpha\). Indeed every chamber is its unique translate by an element
of \(R(F)\), so its image is forced by conjugation; every object belongs
to a chamber. Conversely every signed basis permutation acts by (F6),
fixes the standard chamber, and normalizes the natural group. The
chamber-regular subgroup has trivial intersection with that stabilizer,
and every normalizer element has the unique factorization already
described. This proves (F8). \(\square\)

In particular, the geometry gives a weak representation of the specified
pair \[
(F,\ F\rtimes W_I)
\] when the inner group is written first as in the source's pair
notation. This is a positive statement about compatible injections into
\((K,C)\); it is not a representation of the full abstract group \(F\)
in the source's special sense.

\textbf{Corollary R2.NaturalObstruction.} For \(r\ge2\), there is no
injective homomorphism \[
\operatorname{Aut}F\longrightarrow C(\Gamma_I)
\] whose restriction to \(\operatorname{Inn}F\cong F\) is the natural
embedding \(R\).

\textbf{Proof.} The inner subgroup is normal, so the image of such a
homomorphism would lie in \(N_C(R(F))\). Injectivity, and equality of
the image of the inner subgroup with \(R(F)\), would induce an injection
\[
\operatorname{Out}F\hookrightarrow N_C(R(F))/R(F)\cong W_I.
\] The right side is finite. The left side is infinite: for distinct
basis elements \(s_1,s_2\), the automorphisms \(s_1\mapsto s_1s_2^m\),
fixing the other basis elements, induce distinct transformations of the
abelianization \(\mathbb Z^r\), whereas inner automorphisms act
trivially there. This is a contradiction. \(\square\)

The corollary fixes the natural embedding used in the construction. It
does not rule out every conceivable weak embedding with a different
image of \(\operatorname{Inn}F\). The source's finite-rank definition
separately rules out an \emph{exact} representation of a nonabelian free
group: its outer quotient would have to embed into the finite palette
permutation group. That observation is already made in {[}L126, §7{]}.

\section{Weak representations from panel permutations}
The natural-normalizer obstruction fixes the prescribed chamber-regular
inner action. It does not obstruct other compatible embeddings in the
source's definition of a weak representation. The following construction
settles that distinction without changing the incidence geometry.

\begin{theorem}[An unrestricted panel-permutation subgroup]
\label{UR.PanelSymmetric}
For a finite free basis of size $r\ge2$, the type-preserving group of
$\Gamma_I$ contains $\Sym(\mathbb Z\setminus\{0\})$. Consequently
it contains every countable group.
\end{theorem}
\begin{proof}
Fix $i\in I$. Write elements of $F$ in their unique reduced syllable
form $s_{j_1}^{n_1}\cdots s_{j_k}^{n_k}$, with nonzero integer
exponents and consecutive indices different. Extend any
$\pi\in\Sym(\mathbb Z\setminus\{0\})$ by $\pi(0)=0$. Define
$\Phi_\pi:F\to F$ as follows: if the final syllable is $s_i^n$,
replace it by $s_i^{\pi(n)}$; otherwise leave the word unchanged.
The empty word is fixed. Nonzero exponents remain nonzero, so reduced
syllable forms stay reduced. It follows directly that
$\Phi_\pi\Phi_\rho=\Phi_{\pi\rho}$ and
$\Phi_\pi^{-1}=\Phi_{\pi^{-1}}$.

For $j\in I$, call $g,h$ $j$-adjacent when
$h=s_j^m g$ for some integer $m$; these are precisely the chambers in
one $j$-panel. The map $\Phi_\pi$ preserves this relation in both
directions. On the panel $\langle s_i\rangle$ this is immediate: it
is an arbitrary permutation of that panel fixing its identity chamber.
Outside that case, multiplication on the left changes or cancels only
the first syllable. It leaves the final $i$-syllable unchanged, if one
is present; even when a two-syllable word loses its first syllable,
both orders of the two operations give $s_i^{\pi(n)}$. If no final
$i$-syllable is present the words remain fixed, except for the already
handled panel at the identity. Thus adjacent images are again in the
same $j$-panel. Applying the inverse proves the converse.

The chambers in an object $F_jg$ are exactly those reachable from $g$
using panels of labels other than $j$, because $F_j$ is generated by
the corresponding cyclic factors. Therefore $\Phi_\pi$ permutes the
right cosets of every $F_j$. Define its object action by
$F_jg\mapsto F_j\Phi_\pi(g)$. The preceding observation proves
well-definedness and bijectivity. Two objects are incident exactly when
their sets of chambers intersect, so this action preserves incidence
and types. Distinct permutations $\pi$ give distinct object actions: the type
$i$ cosets $F_is_i^n$ are distinct, since
$F_i\cap\langle s_i\rangle=1$, and their images detect $\pi(n)$.
This proves the symmetric-group embedding.

Finally a countable group embeds faithfully in the symmetric group of
its underlying set by its regular action. Identify that set with a
subset of $\mathbb Z\setminus\{0\}$ and fix the complement to obtain
the asserted embedding.
\end{proof}

\begin{corollary}[Weak representation of a finite-rank free group]
\label{UR.WeakFree}
The geometry $\Gamma_I$ is a weak incidence representation of the
abstract free group $F$ in the source's precise sense. Both compatible
images may be taken inside its type-preserving group.
\end{corollary}
\begin{proof}
An automorphism of $F$ is determined by the images of its finitely many
basis elements. The set $F^r$ is countable, so $\Aut(F)$ is countable.
Embed $\Aut(F)$ in $K(\Gamma_I)$ using
Theorem~\ref{UR.PanelSymmetric}, and restrict this injection to
$\operatorname{Inn}(F)$. Composing the first map with
$K(\Gamma_I)\hookrightarrow C(\Gamma_I)$ gives the required commuting
pair of injections.
\end{proof}
This is an existence construction with a different inner image. It
neither extends the natural chamber-regular action nor contradicts
Theorem R2.Normalizer and its natural-embedding obstruction. No quotient
identification is inferred from these weak injections.

\section{Chamber stabilizers and the splitting issue}
A geometry is \emph{thin} when every panel belongs to exactly two chambers.
The chamber graph joins chambers differing in one object. This graph need not
be connected merely from the definition that every flag extends to a chamber.

\begin{proposition}[The connected case]\label{UR.ThinSplit}
Suppose a thin geometry has a connected chamber graph and its type-preserving
automorphism group $K$ is chamber-transitive. Write $C$ for its full correlation
group. Then $K$ acts regularly on chambers. For any chamber $D$, its setwise
stabilizer $C_D$ is a complement to $K$ in $C$:
\[C=K\rtimes C_D,\qquad C_D\cong C/K.\]
\end{proposition}
\begin{proof}
An element of $K$ fixing $D$ fixes its objects individually. It fixes each
adjacent chamber: the common panel is fixed, and of its two chambers one is
already fixed, so the other is fixed as well. Connectedness propagates this
to every chamber. Every object lies in a chamber; hence the element is the
identity. Chamber transitivity now gives regularity. Given $c\in C$, choose
$k\in K$ with $kD=cD$. Then $k^{-1}c\in C_D$. Thus $C=KC_D$, while
$K\cap C_D=1$. Normality of the palette kernel $K$ proves the semidirect
product and the quotient identification.
\end{proof}

\begin{proposition}[Thin doubling of a bipartite graph]\label{UR.ThinDoubling}
Let $B=(P\sqcup Q,E)$ be a finite connected simple bipartite graph with
$|P|=|Q|\ne|E|$. Assume that its part-preserving automorphism group
$H^+$ is transitive on $E$ and that some automorphism exchanges its parts.
Write $H=\Aut(B)$, so $H/H^+\cong C_2$. There is a finite thin,
flag-transitive rank-three geometry $\Delta(B)$ with
\[
 C(\Delta(B))\cong C_2^{P\sqcup Q\sqcup E}\rtimes H,
 \qquad
 K(\Delta(B))\cong C_2^{P\sqcup Q\sqcup E}\rtimes H^+.
\]
If the exchanging coset $H\setminus H^+$ contains no involution, the
first group does not split over the second. Its chamber graph has one
component for each edge of $B$.
\end{proposition}
\begin{proof}
For each $v\in P\sqcup Q\sqcup E$ take two objects $(v,0),(v,1)$,
with the three displayed sets as their types. Incidence ignores the
second coordinate: a $P$-object and a $Q$-object are incident precisely
when their underlying vertices form an edge of $B$; an edge-object is
incident with precisely the objects over its two endpoints. Distinct objects of one type are nonincident; we use the source's
irreflexive incidence convention.
A chamber is exactly an edge together with independent choices of one
of the two objects over each endpoint and that edge. Every incident
pair extends to such a chamber, as does every single object since the
graph is connected and has an edge. Thus every flag extends to a
chamber. For a panel, the underlying missing vertex or edge is unique,
and its two copies give exactly two completions. This proves thinness.

The independent flips of all pairs form
$N=C_2^{P\sqcup Q\sqcup E}$. The group $H$ lifts by leaving second
coordinates unchanged. The subgroup $N\rtimes H^+$ is transitive on
chambers, using edge transitivity followed by the three endpoint/edge
flips. It is therefore transitive on flags of each type.

Every correlation fixes the edge type: its size $2|E|$ differs from
both other type sizes. It either fixes or exchanges the other two
types. Within each type the two copies of an underlying object are
exactly the classes with equal neighborhoods in the other types.
For distinct vertices of $P$ or $Q$, their sets of incident edge-objects
are disjoint and nonempty, so their neighborhoods differ. Distinct
edge-objects have distinct pairs of endpoints because $B$ is simple.
Consequently these pairs are intrinsically determined, and every
correlation descends to an automorphism of $B$, possibly exchanging
its parts. Its action on edges is forced by its action on vertices.
The restriction kernel consists exactly of the independent flips $N$.
The explicit lifts split this restriction and prove both group formulas.

The quotient by the type-preserving subgroup is $C_2$. A complement
would contain an involution exchanging the two vertex types. Its
projection to $H$ would be an involution in $H\setminus H^+$,
contrary to the stated hypothesis. Finally adjacent chambers have the
same underlying edge; the eight chambers over each edge form the
connected three-dimensional cube under changes of a single choice.
These are precisely the chamber components.
\end{proof}

The second application supplies a finite graph satisfying the nonsplitting
hypothesis, with a complete elementary certificate of its full automorphism
group. Thus the unqualified source observation is false, whereas the
connected-chamber statement above holds by its direct proof.

The following inherited verification discussion concerns the two programs
and restriction module for the first application. The second application's
separate scope record also includes the full optional finite/free-word checks,
the required graph certificate with its complete data, and the inherited
basis-permutation and shadow-kernel formal declarations. The certificate is
needed only for the specific nonsplit graph, not for the generic doubling
argument.
\section{Attribution and verification}
The two constructions, their symmetry assertions and their use in the finite
normal-pair theorem are due to D\'iaz Ramos, Molinier and Viruel. The
restriction-and-extension criterion is elementary classical reasoning; its
explicit separation from recognition and rigidity is an organization of the
proofs, not a priority claim. Lemma~\ref{UR.Components} isolates the component
argument also used in the recursive application. No external graph
representation theorem is needed to make the proofs above complete.

The accompanying \texttt{RecognizableRestriction.lean} constructs the
restriction homomorphism from its actual action on an embedded vertex set,
proves its kernel description, derives injectivity from pointwise rigidity and
surjectivity from existence of extensions, constructs the group equivalence,
and transports a distinguished subgroup. It does not formalize the incidence
or degree arguments proving those hypotheses for any gadget. The complete
chamber, cleaning and component arguments above remain written proofs.

The two complete GAP programs in the separate appendix provide bounded
checks, including exceptional graphs and small normal pairs. Their computations
do not replace any quantified theorem above. GAP and GRAPE's exact finite
automorphism algorithms remain part of the trust boundary of those checks;
Lean uses its kernel and the pinned mathlib infrastructure. Neither the number
of programs nor the number of formal declarations measures mathematical
compression.

\begin{thebibliography}{9}
\bibitem{DMV} A. D\'iaz Ramos, R. Molinier and A. Viruel,
\emph{Every finite group is represented by a finite incidence geometry},
arXiv:2511.02410v2, 16 December 2025.
\end{thebibliography}

\end{document}
